\documentclass[tikz,border=1.5mm]{standalone}
\input{src/_preamble.tex}
\begin{document}
\begin{tikzpicture}[x=1mm,y=1mm]
  % 每个面板：4 行（左前 右前 左后 右后），横轴为一个步态周期（0~1），宽 46 mm
  % \gait{x}{y}{标题}{LF起,止;...}
  \newcommand{\panel}[5]{%
    \begin{scope}[shift={(#1,#2)}]
      \node[font=\scriptsize\bfseries, anchor=west] at (0,23) {#3};
      \node[font=\tiny, text=mid, anchor=east] at (-1,15) {左前};
      \node[font=\tiny, text=mid, anchor=east] at (-1,10) {右前};
      \node[font=\tiny, text=mid, anchor=east] at (-1,5) {左后};
      \node[font=\tiny, text=mid, anchor=east] at (-1,0) {右后};
      \foreach \yy in {0,5,10,15} \fill[black!7] (0,\yy-1.6) rectangle (46,\yy+1.6);
      \foreach \yy/\a/\b in {#4} \fill[ink] (\a*46,\yy-1.6) rectangle (\b*46,\yy+1.6);
      \draw[mid, line width=0.3pt] (0,-3) -- (46,-3);
      \node[font=\tiny, text=mid, anchor=north] at (0,-3) {0};
      \node[font=\tiny, text=mid, anchor=north] at (46,-3) {一个周期};
      \node[font=\tiny, text=mid, anchor=west, align=left] at (0,-8.5) {#5};
    \end{scope}}
  % 慢走：每条腿着地 75%，依次抬腿（左后→左前→右后→右前）
  \panel{0}{40}{慢走}{15/0/0.25, 15/0.5/1.0, 10/0/0.75, 5/0.25/1.0, 0/0/0.5, 0/0.75/1.0}{任何时刻至少三只脚着地}
  % 对角小跑：左前+右后、右前+左后
  \panel{62}{40}{对角小跑}{15/0/0.5, 0/0/0.5, 10/0.5/1.0, 5/0.5/1.0}{对角两条腿一起落地}
  % 同侧溜蹄：左前+左后、右前+右后
  \panel{0}{0}{同侧溜蹄}{15/0/0.5, 5/0/0.5, 10/0.5/1.0, 0/0.5/1.0}{同一侧两条腿一起落地}
  % 跳跑：前两条、后两条
  \panel{62}{0}{跳跑}{15/0.5/0.85, 10/0.5/0.85, 5/0/0.35, 0/0/0.35}{前腿一起、后腿一起，中间有腾空}
  % 图例
  \fill[ink] (0,-17) rectangle (5,-14); \node[font=\scriptsize, anchor=west] at (6,-15.5) {着地};
  \fill[black!7] (24,-17) rectangle (29,-14); \node[font=\scriptsize, anchor=west] at (30,-15.5) {在空中摆动};
\end{tikzpicture}
\end{document}
